\section{问题背景与重述}
\subsection{问题背景}
塔式太阳能光热发电通过定日镜阵列将太阳辐射反射至吸收塔顶部的集热器，使导热介质升温并与储热、发电系统衔接。定日镜既是采集太阳能的基本单元，也是镜场占地和设备数量的主要来源。提高镜场的面积利用效率，有助于在满足额定产热要求的同时减少反射镜材料用量，是电站前期设计中的关键任务\cite{problem,zhang}。

定日镜场的光学收益并不简单正比于镜面面积。一方面，太阳高度与方位随季节和时刻变化，镜面需要持续调整姿态；另一方面，定日镜之间会在入射与反射路径上发生阴影和遮挡。太阳具有有限角直径，平面镜反射的光束也具有空间展宽，部分能量不能落到有限尺寸的集热器上。这些损失同时受塔址、镜位、镜面尺寸和安装高度影响，因此应在完整几何下计算，再对设计变量进行优化。

本文研究的场址位于东经 $98.5^\circ$、北纬 $39.4^\circ$、海拔 $3000\,\mathrm m$，可建设区域是半径 $350\,\mathrm m$ 的水平圆盘。集热器中心高 $80\,\mathrm m$，直径 $7\,\mathrm m$、高度 $8\,\mathrm m$，塔周 $100\,\mathrm m$ 范围为禁布区。镜面上下边保持水平，镜宽与镜高均在 $2$ 至 $8\,\mathrm m$ 之间，安装高度为 $2$ 至 $6\,\mathrm m$；镜面旋转不得触地，相邻底座还需满足维护通道的距离要求。

\subsection{问题重述}
建立以场地圆心为原点、正东为 $x$ 轴、正北为 $y$ 轴、竖直向上为 $z$ 轴的坐标系。按照题目规定，每月 $21$ 日的 $9{:}00$、$10{:}30$、$12{:}00$、$13{:}30$ 和 $15{:}00$ 为计算时刻，全年共 $60$ 个时刻。需要解决以下三个相互衔接的问题。

\textbf{问题一：给定镜场的光学评价。}吸收塔位于圆心，镜面统一为 $6\,\mathrm m\times6\,\mathrm m$，安装高度统一为 $4\,\mathrm m$，镜心位置由附件给出。计算各月和全年的平均光学效率、输出热功率及单位镜面面积输出热功率。

\textbf{问题二：统一镜型的布局设计。}要求所有镜面尺寸和安装高度相同，同时设计吸收塔位置、镜面尺寸、安装高度、镜数和镜位。在年均热功率不低于 $60\,\mathrm{MW}$ 时，尽量提高单位镜面面积年均热功率。

\textbf{问题三：异构镜型的布局设计。}允许不同定日镜采用不同尺寸和安装高度，在相同额定容量要求下重新分配镜面面积与空间位置，给出可行方案并分析相对统一镜型的改进。

\pagepart
\section{问题分析}
\textbf{问题一的核心是把相关损失放在同一光路中计算。}由附件确定镜心之后，太阳方向和镜面法向即可在各时刻唯一确定。余弦效率、大气透射率可直接计算，阴影、遮挡和截断则需要光线与有限几何体求交。本文采用有限太阳圆盘与矩形镜面联合采样，先判定入射方向可见性，再判定反射路径和接收表面。这样既保留太阳光锥角，也能避免分别扣除重叠损失。最后先逐镜累计功率，再按题目规定的时刻等权平均，得到定场评价结果。

\textbf{问题二的核心是协调容量、镜型与排列密度。}扩大镜面能增加采光面积，却会提高底座中心距下限，并可能加重遮挡与截断损失；塔址偏南可能改善部分区域的余弦效率，但也会改变可用场地和反射距离。直接优化数千个坐标的成本很高，本文使用以塔为中心的分区交错同心圆环表示布局，以弦距约束确定紧凑环距。通过尺寸、塔址与区段圈数的组合搜索生成候选，再在容量余量允许时删除低单位面积收益镜面，完成面积压缩。

\textbf{问题三的核心是异构面积分配与遮挡外部效应之间的平衡。}每面镜的面积收益不同，因此既不应要求远镜更大，也不应强制各镜输出一致。本文先以径向分层降低安装高度变量的维数，再在固定邻居条件下计算各镜的离散规格收益，以拉格朗日面积价格统一协调“保留、缩小、删除”等选项。代理模型用于快速提出方案，全场耦合追迹用于判定实际容量和目标值；两者共同控制计算成本与结果可信度。

\begin{figure}[H]\centering
\begin{tikzpicture}[>=Latex,node distance=9mm,
 box/.style={draw=navy,very thick,fill=navy!5,rounded corners=2pt,align=center,text width=112mm,minimum height=10mm,font=\small},
 stagebox/.style={draw=orange,thick,fill=orange!7,align=center,text width=33mm,minimum height=22mm,font=\small}]
\node[box] (data) {场址、镜心与设备约束；全年 60 个规定时刻};
\node[box,below=of data] (trace) {太阳位置与镜面姿态 $\rightarrow$ 有限光锥采样 $\rightarrow$ 完整光路评价};
\node[stagebox,below=12mm of trace] (two) {问题二\\交错圆环、组合搜索\\容量约束下删镜};
\node[stagebox,left=6mm of two] (one) {问题一\\给定镜场逐镜计算\\月均与年均汇总};
\node[stagebox,right=6mm of two] (three) {问题三\\分层高度、面积价格\\离散规格与耦合复算};
\node[box,below=13mm of two] (check) {几何检查、功率恒等式、采样收敛与参数敏感性};
\draw[->,thick] (data)--(trace);
\draw[->,thick] (trace.south)--(two.north);
\draw[->,thick] (trace.south west)+(18mm,0)--(one.north);
\draw[->,thick] (trace.south east)+(-18mm,0)--(three.north);
\draw[->,thick] (one.south)--($(check.north west)+(18mm,0)$);
\draw[->,thick] (two)--(check);
\draw[->,thick] (three.south)--($(check.north east)+(-18mm,0)$);
\end{tikzpicture}
\caption{三问建模与求解的总体流程}
\end{figure}

\pagepart
\section{模型假设}
为使几何、辐照度和设计约束闭合，采用以下假设。题目直接给出的条件与额外采用的建模约定分别说明。
\begin{enumerate}[label=\textbf{假设\arabic*：},leftmargin=4.8em]
\item \textbf{场地水平且无遮挡地形。}地面取 $z=0$，不引入山体、建筑物或地形坡度。镜间阴影、吸收塔和集热器本体仍显式参与计算。
\item \textbf{规定时刻等权。}采用题目附录的太阳时角公式，以 $3$ 月 $21$ 日为春分零日，不另作北京时间换算。各月均取五个时刻，年均取十二个月的等权平均，不进行月份天数和日照时长加权。
\item \textbf{镜面为理想平面反射镜。}镜面反射率取 $0.92$，跟踪系统使太阳中心光线经镜心反射后指向集热器中心；暂不计跟踪误差、面形误差、积灰和镜面老化。
\item \textbf{太阳采用均匀圆盘光源。}太阳半角取 $4.65\,\mathrm{mrad}$，采样方向由圆盘光锥产生。辐照能量按太阳中心方向的余弦因子加权，这是与求解程序一致的小角度近似。
\item \textbf{集热器侧面有效受光。}集热器视为半径 $3.5\,\mathrm m$、高度区间 $[76,84]\,\mathrm m$ 的闭圆柱；侧面首次命中计为接收，端面作为不透明边界但不计入有效吸收面积。
\item \textbf{塔身尺寸采用显式补充假设。}塔身视为半径 $3.5\,\mathrm m$、高度区间 $[0,76]\,\mathrm m$ 的不透明圆柱。题目未给塔身半径，因此在第六部分单独检验这一设定的影响。
\item \textbf{场界和禁区约束作用于镜心。}题目将定日镜位置定义为中心位置，本文据此要求镜心位于场地圆盘内且在塔周禁布区外。不同镜宽的相邻镜对，以较大镜宽加 $5\,\mathrm m$ 作为底座间距下限。
\item \textbf{安装高度满足全旋转离地条件。}采用 $z_i-h_i/2>0$ 保证任意俯仰下镜面不触地。第二问统一镜型搜索留出 $0.05\,\mathrm m$ 的几何余量，第三问保留相同或更高的安装中心。
\item \textbf{研究对象为接收热功率。}不把光学接收功率继续换算为电功率，不引入吸热器热损、储热损失或热电转换效率；所得结论适用于题目定义的镜场光学评价。
\end{enumerate}
这些假设保证三问使用一致的评价基础。特别地，紧凑布置不代表取消维护间距，分层抬高也不代表免除全场遮挡计算。模型误差与随机采样误差的来源不同，不能用增加光线数消除物理假设偏差。

\pagepart
\section{符号说明}
\begin{table}[H]\centering\small
\caption{主要符号、说明与单位}
\begin{tabular}{p{30mm}p{95mm}p{20mm}}\toprule
符号 & 说明 & 单位\\\midrule
$\mathcal T$，$N$ & 规定时刻集合、定日镜数量 & ---，面\\
$\vct c_i=(x_i,y_i,z_i)$ & 第 $i$ 面镜的中心坐标，$z_i$ 为安装高度 & m\\
$w_i,h_i,A_i$ & 镜面宽、高、面积，$A_i=w_ih_i$ & m，m$^2$\\
$\vct T=(x_T,y_T,H_T)$ & 集热器中心；$H_T=80$ & m\\
$R_f,R_e,R_c$ & 场地半径、禁区半径、集热器半径 & m\\
$\varphi,\delta,\omega$ & 场址纬度、太阳赤纬、太阳时角 & rad\\
$D,ST$ & 相对春分的日数、当地太阳时 & d，h\\
$\vct s,\vct s'$ & 太阳中心方向、采样太阳方向单位向量 & ---\\
$\vct r_i,\vct n_i$ & 镜心至接收中心方向、镜面法向单位向量 & ---\\
$\vct u_i,\vct v_i$ & 镜面水平宽度、高度方向单位向量 & ---\\
$d_i$ & 镜心到集热器中心的距离 & m\\
$\alpha_\odot$ & 太阳光锥半角 & rad\\
$K$，$\phi_b(k)$ & 每镜每时刻光线数、基数 $b$ 的逆根函数 & ---\\
$S_{itk},B_{itk},H_{itk}$ & 入射可见、反射可见、接收命中指示量 & ---\\
$\eta_{\rm cos},\eta_{\rm at}$ & 余弦效率、大气透射率 & ---\\
$\eta_{\rm sb},\eta_{\rm tr}$ & 阴影遮挡效率、条件截断效率 & ---\\
$\eta_{\rm ref},\eta_i$ & 镜面反射率、单镜总光学效率 & ---\\
$\DNI_t$ & 时刻 $t$ 的法向直接辐照度 & kW/m$^2$\\
$P_t,\mean P,P_*$ & 瞬时全场功率、年均功率、额定功率 & kW\\
$A_\Sigma,J$ & 镜面总面积、单位面积年均热功率 & m$^2$，kW/m$^2$\\
$p,r_k,n,g$ & 布置节距、环半径、名义镜数、区段圈数 & m，面，圈\\
$\Delta r_k,\psi$ & 环间径向增量、无量纲角相位 & m，---\\
$z_\ell,\mathcal I_\ell$ & 第 $\ell$ 个径向区间的安装高度与镜子集合 & m，---\\
$P_{ij}^{\rm pr},A_{ij}$ & 镜 $i$ 采用选项 $j$ 的代理功率与面积 & kW，m$^2$\\
$\lambda,x_{ij}$ & 面积价格、离散选项决策变量 & kW/m$^2$，---\\
\bottomrule\end{tabular}
\end{table}
全文推导中功率统一以 kW 计算，结果表中全场功率转换为 MW；单位面积热功率始终以 kW/m$^2$ 表示。上横线代表对规定时刻平均，下标 $\ell$ 表示径向高度区间，不能与底座圆环序号 $k$ 混用。表中的尺寸保留有限小数，几何约束使用完整精度数据检查。
\pagepart
